\documentclass[10pt,a4paper]{amsart}
\usepackage{amsmath,amssymb,amsthm,graphicx}
% Retained preamble for the Stacks Project chapter snapshot in this directory.
% The chapter's own preamble.tex is a copy of the shared Stacks Project preamble; the block
% below is sliced MECHANICALLY, line for line, from it. Nothing was added, removed or
% redefined; the wrapper class, the xr-hyper/hyperref wiring and the externaldocument list
% are omitted because the assembler supplies the document class.
%
% lines taken from preamble.tex: 3, 5, 9, 11, 12, 156, 157, 158, ... (51 lines total)
%
% For commutative diagrams we use Xy-pic
\usepackage[all]{xy}

% We use 2cell for 2-commutative diagrams.
\xyoption{2cell}
\UseAllTwocells

\newenvironment{reference}{\comment}{\endcomment}
\newenvironment{slogan}{\comment}{\endcomment}
\newenvironment{history}{\comment}{\endcomment}
\theoremstyle{plain}
\newtheorem{theorem}[subsection]{Theorem}
\newtheorem{proposition}[subsection]{Proposition}
\newtheorem{lemma}[subsection]{Lemma}
\theoremstyle{definition}
\newtheorem{definition}[subsection]{Definition}
\newtheorem{example}[subsection]{Example}
\newtheorem{exercise}[subsection]{Exercise}
\newtheorem{situation}[subsection]{Situation}
\theoremstyle{remark}
\newtheorem{remark}[subsection]{Remark}
\newtheorem{remarks}[subsection]{Remarks}
\numberwithin{equation}{subsection}
\def\lim{\mathop{\mathrm{lim}}\nolimits}
\def\colim{\mathop{\mathrm{colim}}\nolimits}
\def\Spec{\mathop{\mathrm{Spec}}}
\def\Hom{\mathop{\mathrm{Hom}}\nolimits}
\def\Ext{\mathop{\mathrm{Ext}}\nolimits}
\def\SheafHom{\mathop{\mathcal{H}\!\mathit{om}}\nolimits}
\def\SheafExt{\mathop{\mathcal{E}\!\mathit{xt}}\nolimits}
\def\Sch{\mathit{Sch}}
\def\Mor{\mathop{\mathrm{Mor}}\nolimits}
\def\Ob{\mathop{\mathrm{Ob}}\nolimits}
\def\Sh{\mathop{\mathit{Sh}}\nolimits}
\def\NL{\mathop{N\!L}\nolimits}
\def\CH{\mathop{\mathrm{CH}}\nolimits}
\def\proetale{{pro\text{-}\acute{e}tale}}
\def\etale{{\acute{e}tale}}
\def\QCoh{\mathit{QCoh}}
\def\Ker{\mathop{\mathrm{Ker}}}
\def\Im{\mathop{\mathrm{Im}}}
\def\Coker{\mathop{\mathrm{Coker}}}
\def\Coim{\mathop{\mathrm{Coim}}}
\DeclareMathSymbol{\boxtimes}{\mathbin}{AMSa}{"02}
\def\QCohstack{\mathcal{QC}\!\mathit{oh}}
\def\Cohstack{\mathcal{C}\!\mathit{oh}}
\def\Spacesstack{\mathcal{S}\!\mathit{paces}}
\def\Quotfunctor{\mathrm{Quot}}
\def\Hilbfunctor{\mathrm{Hilb}}
\def\Curvesstack{\mathcal{C}\!\mathit{urves}}
\def\Polarizedstack{\mathcal{P}\!\mathit{olarized}}
\def\Complexesstack{\mathcal{C}\!\mathit{omplexes}}
\def\Pic{\mathop{\mathrm{Pic}}\nolimits}
\def\Picardstack{\mathcal{P}\!\mathit{ic}}
\def\Picardfunctor{\mathrm{Pic}}
\def\Deformationcategory{\mathcal{D}\!\mathit{ef}}
\usepackage{cleveref}
\crefname{originalsection}{original section}{original sections}
\Crefname{originalsection}{Original section}{Original sections}
\makeatletter
\newlabel{section-loc-qcoh-flat-base-change}{{0GQG}{source}{}{sourcepointer-1}{}}
\newlabel{section-loc-qcoh-flat-base-change@cref}{{[originalsection][0][]0GQG}{[1][0][]source}{}{}{}}
\makeatother
\crefname{originalreference}{original reference}{original references}
\Crefname{originalreference}{Original reference}{Original references}
\makeatletter
\newlabel{algebraic-lemma-stack-presentation}{{04T5}{source}{}{sourcepointer-2}{}}
\newlabel{algebraic-lemma-stack-presentation@cref}{{[originalreference][0][]04T5}{[1][0][]source}{}{}{}}
\makeatother
\crefname{originalreference}{original reference}{original references}
\Crefname{originalreference}{Original reference}{Original references}
\makeatletter
\newlabel{stacks-sheaves-proposition-quasi-coherent}{{06WT}{source}{}{sourcepointer-3}{}}
\newlabel{stacks-sheaves-proposition-quasi-coherent@cref}{{[originalreference][0][]06WT}{[1][0][]source}{}{}{}}
\makeatother
\crefname{originalreference}{original reference}{original references}
\Crefname{originalreference}{Original reference}{Original references}
\makeatletter
\newlabel{spaces-groupoids-definition-groupoid-module}{{0441}{source}{}{sourcepointer-4}{}}
\newlabel{spaces-groupoids-definition-groupoid-module@cref}{{[originalreference][0][]0441}{[1][0][]source}{}{}{}}
\makeatother
\crefname{originalreference}{original reference}{original references}
\Crefname{originalreference}{Original reference}{Original references}
\makeatletter
\newlabel{stacks-sheaves-lemma-map-from-quasi-coherent}{{076S}{source}{}{sourcepointer-5}{}}
\newlabel{stacks-sheaves-lemma-map-from-quasi-coherent@cref}{{[originalreference][0][]076S}{[1][0][]source}{}{}{}}
\makeatother
\title{The right adjoint Q of the inclusion QCoh into LQCoh\textasciicircum{}\{fbc\}}
\author{The Stacks Project Authors}
\date{stacks-project:stacks-cohomology.tex:local-snapshot-20261012; GNU FDL 1.2 or later}
\begin{document}
\maketitle
\noindent\textit{Editorial note.} All delivered mathematical text is the author's own, copied verbatim from the pinned Stacks Project chapter file stacks-cohomology.tex; no statement or proof step was rewritten, reordered or removed.\par\medskip
\noindent\textit{Editorial note.} The author omits the routine verification of the cocycle condition for the descent datum alpha; the construction itself is complete.\par\medskip
\noindent\textit{Allowed standard prerequisites (editorial).}\begin{itemize}
\item Stack presentation by a smooth groupoid in spaces. Algebraic Stacks, Lemma algebraic-lemma-stack-presentation; Sheaves on Stacks, Proposition stacks-sheaves-proposition-quasi-coherent. Presentation and equivalence.
\item Groupoid modules descent datum. Groupoids in Spaces, Definition spaces-groupoids-definition-groupoid-module. Cocycle condition (omitted routine check).
\end{itemize}\medskip
\ifdefined\hypertarget\hypertarget{sourcepointer-1}{}\fi
\noindent\textit{Source reference.} Original section 0GQG: sources/stacks-chow/tags.txt, source lines 20861--20861. Cross-chapter standard prerequisite cited by the author inside the counted span; it is used as a pointer and its own proof is not part of this item.\par\medskip
\ifdefined\hypertarget\hypertarget{sourcepointer-2}{}\fi
\noindent\textit{Source reference.} Original reference 04T5: sources/stacks-chow/tags.txt, source lines 6243--6243. Cross-chapter standard prerequisite cited by the author inside the counted span; it is used as a pointer and its own proof is not part of this item.\par\medskip
\ifdefined\hypertarget\hypertarget{sourcepointer-3}{}\fi
\noindent\textit{Source reference.} Original reference 06WT: sources/stacks-chow/tags.txt, source lines 8827--8827. Cross-chapter standard prerequisite cited by the author inside the counted span; it is used as a pointer and its own proof is not part of this item.\par\medskip
\ifdefined\hypertarget\hypertarget{sourcepointer-4}{}\fi
\noindent\textit{Source reference.} Original reference 0441: sources/stacks-chow/tags.txt, source lines 5384--5384. Cross-chapter standard prerequisite cited by the author inside the counted span; it is used as a pointer and its own proof is not part of this item.\par\medskip
\ifdefined\hypertarget\hypertarget{sourcepointer-5}{}\fi
\noindent\textit{Source reference.} Original reference 076S: sources/stacks-chow/tags.txt, source lines 9177--9177. Cross-chapter standard prerequisite cited by the author inside the counted span; it is used as a pointer and its own proof is not part of this item.\par\medskip
\begin{lemma}
\label{lemma-adjoint}
Let $\mathcal{X}$ be an algebraic stack. Let
$\textit{LQCoh}^{fbc}(\mathcal{O}_\mathcal{X})$
be the category of locally quasi-coherent modules with the
flat base change property, see
Section \ref{section-loc-qcoh-flat-base-change}.
The inclusion functor
$i : \QCoh(\mathcal{O}_\mathcal{X}) \to
\textit{LQCoh}^{fbc}(\mathcal{O}_\mathcal{X})$
has a right adjoint
$$
Q : \textit{LQCoh}^{fbc}(\mathcal{O}_\mathcal{X}) \to
\QCoh(\mathcal{O}_\mathcal{X})
$$
such that $Q \circ i$ is the identity functor.
\end{lemma}

\begin{proof}
Choose a scheme $U$ and a surjective smooth morphism $f : U \to \mathcal{X}$.
Set $R = U \times_\mathcal{X} U$ so that we obtain a smooth groupoid
$(U, R, s, t, c)$ in algebraic spaces with the property that
$\mathcal{X} = [U/R]$, see
Algebraic Stacks, Lemma \ref{algebraic-lemma-stack-presentation}.
We may and do replace $\mathcal{X}$ by $[U/R]$. By
Sheaves on Stacks, Proposition \ref{stacks-sheaves-proposition-quasi-coherent}
there is an equivalence
$$
q_1 :
\QCoh(U, R, s, t, c)
\longrightarrow
\QCoh(\mathcal{O}_\mathcal{X})
$$
Let us construct a functor
$$
q_2 :
\textit{LQCoh}^{fbc}(\mathcal{O}_\mathcal{X})
\longrightarrow
\QCoh(U, R, s, t, c)
$$
by the following rule: if $\mathcal{F}$ is an object of
$\textit{LQCoh}^{fbc}(\mathcal{O}_\mathcal{X})$ then we set
$$
q_2(\mathcal{F}) = (f^*\mathcal{F}|_{U_\etale}, \alpha)
$$
where $\alpha$ is the isomorphism
$$
t_{small}^*(f^*\mathcal{F}|_{U_\etale})
\to
t^*f^*\mathcal{F}|_{R_\etale} \to
s^*f^*\mathcal{F}|_{R_\etale} \to
s_{small}^*(f^*\mathcal{F}|_{U_\etale})
$$
where the outer two morphisms are the comparison maps. Note that
$q_2(\mathcal{F})$ is quasi-coherent precisely because $\mathcal{F}$ is
locally quasi-coherent and that we used (and needed)
the flat base change property in the construction of
the descent datum $\alpha$. We omit the
verification that the cocycle condition (see
Groupoids in Spaces, Definition
\ref{spaces-groupoids-definition-groupoid-module})
holds. Looking at the proof of
Sheaves on Stacks, Proposition \ref{stacks-sheaves-proposition-quasi-coherent}
we see that $q_2 \circ i$ is the quasi-inverse to $q_1$.
We define $Q = q_1 \circ q_2$.
Let $\mathcal{F}$ be an object of
$\textit{LQCoh}^{fbc}(\mathcal{O}_\mathcal{X})$ and
let $\mathcal{G}$ be an object of $\QCoh(\mathcal{O}_\mathcal{X})$.
We have
\begin{align*}
\Mor_{\textit{LQCoh}^{fbc}(\mathcal{O}_\mathcal{X})}
(i(\mathcal{G}), \mathcal{F})
& =
\Mor_{\QCoh(U, R, s, t, c)}(q_2(i(\mathcal{G})), q_2(\mathcal{F})) \\
& =
\Mor_{\QCoh(\mathcal{O}_\mathcal{X})}(\mathcal{G}, Q(\mathcal{F}))
\end{align*}
where the first equality is
Sheaves on Stacks, Lemma \ref{stacks-sheaves-lemma-map-from-quasi-coherent}
and the second equality holds because $q_1 \circ i$ and $q_2$ are quasi-inverse
equivalences of categories. The assertion $Q \circ i \cong \text{id}$
is a formal consequence of the fact that $i$ is fully faithful.
\end{proof}

\end{document}
